\begin{frame}{자주 묻는 질문 (14.2): ``한 번''이 진짜 한 번인가}
  \begin{itemize}
    \item \textbf{Q. 4,096개 환경이 \textbf{진짜} 동시에인가?}
      구조적으로는 \textbf{``한 번''} --- 호출은 한 번, 그 안에서
      수천 개 코어가 \textbf{동시에} 수행.
    \item 물리적으로: GPU의 SM(Streaming Multiprocessor)이
      \kb{워프}(32개 스레드) 단위로 스케줄링 $\Rightarrow$
      4,096개 환경이 \textbf{완벽히} 동시에 끝나는 건 아님 ---
      대부분이 동시에, 일부는 \textbf{다음 워프 스케줄링}에서
    \item 하지만 호출 오버헤드가 \textbf{일단}만 치러지므로,
      순차 4,096회 호출에 비해 총 시간은 \textbf{수천 분의 1}
  \end{itemize}
  \vskip0.5em
  \begin{block}{핵심}
    ``\textbf{완벽한 동시성}''이 아니라,
    \textbf{``호출 오버헤드를 상수로 묶는다''}는 것이 핵심.
  \end{block}
  \vskip0.4em
  {\footnotesize Q. MuJoCo-XLA vs Isaac Sim? 둘 다 ``정밀 물리를
  GPU/TPU로 병렬화''하지만 \textbf{범위}가 다름: MuJoCo-XLA =
  물리만(LCP 포함, 렌더링 없음, CPU MuJoCo와 \textbf{완전히 같은
  물리}). Isaac Sim = 물리+렌더링+센서(RT 코어 필요, 물리 엔진도
  NVIDIA Warp 기반). ``정밀 물리만 GPU로''면 MuJoCo-XLA,
  ``물리+카메라로 판단하는 정책''이면 Isaac Sim.
  }
\end{frame}
